% Hafta 13 — Dersin gereksinim aileleri
% Derleme: py -3.12 tools/latex/derle.py docs/week-13/assets/h13-11-gereksinim-aileleri.tex
\documentclass[tikz,border=8pt]{standalone}
\usepackage{cen429-cizim}
\begin{document}
\begin{tikzpicture}
  \node[baslik] (b) at (0,0) {Dersin gereksinim aileleri};
  \node[altbaslik] (alt) at (b.south west) {her gereksinim bir aileye, her aile bir varlığa bağlanır};

  % Sütun merkezleri: 2.0, 6.6, 11.2 (cm) — satır merkezleri: -2.2, -4.8, -7.4 (cm)

  % Satır 1 — mor
  \node[morkutu=36mm, minimum height=16mm] at (2.0,-2.2) (g1) {\kb{\kod{CEN429-DR}}\\Beklemede veri};
  \node[morkutu=36mm, minimum height=16mm] at (6.6,-2.2) (g2) {\kb{\kod{CEN429-DT}}\\Aktarımda veri};
  \node[morkutu=36mm, minimum height=16mm] at (11.2,-2.2) (g3) {\kb{\kod{CEN429-DU}}\\Kullanımda veri};

  % Satır 2 — teal / yeşil
  \node[kutu=36mm, minimum height=16mm] at (2.0,-4.8) (g4) {\kb{\kod{CEN429-ID}}\\Kimlik doğrulama\\ve bağlama};
  \node[kutu=36mm, minimum height=16mm] at (6.6,-4.8) (g5) {\kb{\kod{CEN429-AP}}\\Uygulama koruması};
  \node[iyikutu=36mm, minimum height=16mm] at (11.2,-4.8) (g6) {\kb{\kod{CEN429-AS}}\\Varlık koruması};

  % Satır 3 — yeşil / turuncu
  \node[iyikutu=36mm, minimum height=16mm] at (2.0,-7.4) (g7) {\kb{\kod{CEN429-CR}}\\Kripto ve anahtarlar};
  \node[uyarikutu=36mm, minimum height=16mm] at (6.6,-7.4) (g8) {\kb{\kod{CEN429-RP}}\\Raporlama};
  \node[uyarikutu=36mm, minimum height=16mm] at (11.2,-7.4) (g9) {\kb{\kod{CEN429-DV}}\\Geliştirme süreci};

  \node[fit=(g1)(g9), inner sep=0] (izgara) {};
  \node[dip, text width=150mm, below=9mm of izgara.south west, anchor=north west] {%
    Kimlik numarası olmayan gereksinim izlenemez; izlenemeyen gereksinim denetlenemez.};
\end{tikzpicture}
\end{document}
